\documentclass[10pt,a4paper]{amsart}
\usepackage[margin=19mm]{geometry}
\usepackage{amsmath,amssymb,mathtools}
\usepackage[scr=rsfs,cal=euler]{mathalfa}
\usepackage{xfrac}
\usepackage[unicode,hidelinks,bookmarks=false]{hyperref}
\newcommand{\ie}{i.e.\ }
\newcommand{\cf}{cf.\ }
\newcommand{\resp}{respectively\ }
\newcommand{\Subsection}{\subsection}
\providecommand{\nobreakdash}{\nobreak\hbox{-}}
\setlength{\emergencystretch}{2em}
\multlinegap0pt
\usepackage{amssymb}
\newtheorem{lemma}{Lemma}
\newcommand{\Lc}{\mathcal{L}}
\newcommand{\R}{\mathbb{R}}
\newcommand{\T}{\mathbb{T}}
\newcommand{\lan}{\langle}
\newcommand{\ran}{\rangle}
\title{The Hamiltonian formalism for the Calogero--Moser derivative nonlinear Schr\"odinger equation on the torus}
\author{Rowan Killip, Katie Marsden,, Monica Vi\c{s}an}
\date{Source: arXiv:2604.09479v1 (CC BY 4.0)}
\begin{document}
\maketitle

\noindent\emph{Source and licence.} The Hamiltonian formalism for the Calogero--Moser derivative nonlinear Schr\"odinger equation on the torus. Version arXiv:2604.09479v1 (CC BY 4.0), distributed by arXiv under the Creative Commons Attribution 4.0 International licence, http://creativecommons.org/licenses/by/4.0/, as recorded in the arXiv licence metadata for this exact version and shown on the abstract page https://arxiv.org/abs/2604.09479v1. Full licence notice: \texttt{licenses/arxiv-2604.09479-CC-BY-4.0.txt}.\par\smallskip

\noindent\textit{Editorial scope.} Counted unit: the article's lemma H (source0.tex lines 1029--1037). The article's complete original proof of it is delivered with it (source0.tex lines 1039--1146, one \texttt{proof} environment of 108 lines). No mathematics is written, repaired or re-derived: the statement and the proof are the article's own lines, in the article's own notation, and no retained line is substituted or re-flowed.  Retained uncounted as the source-local context the proof needs: lines 282--284.  Nothing was omitted from any retained span, so the package carries no omission record.  No retained line contains a citation, so the wrapper carries no bibliography and no bibliography entry.  No retained line contains a figure environment, an image-inclusion command or drawing code, so no image is delivered and the package declares no asset dependency.  Editorial formatting only, in addition: an AMS \texttt{amsart} preamble that restates the article's own declarations and macros used by the retained lines, namely the environments \texttt{lemma} on separate counters, together with the standard packages those lines need.  The retained environments print in this document as Lemma 1 in order of appearance, whereas the article prints the same items with the numbering of its own full document.  The complete native source of the article as distributed in the arXiv e-print for this exact version is bundled unchanged as \texttt{sources/198210/source0.tex}.
\begin{align}\label{Cid}
C_++C_-=1 \quad \text{on $\R$,\quad but} \quad C_+ f + C_- f = f + \widehat f(0) \quad \text{ on $\T$.}
\end{align}

\begin{lemma}\label{L:P and H}
On $\T$ we have 
\begin{align*}
P(q)&=-\tfrac{1}{2}\|q\|_{\dot{H}^{1/2}}^2\pm \tfrac14\|q\|_{L^4}^4\mp\tfrac1{8\pi}\|q\|_{L^2}^4,\\
H(q)&=\tfrac12\|q'\mp i qC_+(|q|^2)\|_{L^2}^2\pm \tfrac{3}{4\pi} \|q\|_{L^2}^2\|q\|_{\dot{H}^{1/2}}^2-\tfrac{3}{8\pi}\|q\|_{L^2}^2\|q\|_{L^4}^4+\tfrac{1}{16\pi^2}\|q\|_{L^2}^6\notag\\
&= \tfrac12\|q\|_{\dot{H}^{1}}^2\mp\tfrac12\| C_+(|q|^2)\|_{\dot{H}^{1/2}}^2\mp\tfrac14\|q^2\|_{\dot{H}^{1/2}}^2\pm \tfrac1{2\pi}\|q\|_{L^2}^2\|q\|_{\dot{H}^{1/2}}^2\nonumber\\
&\quad+\tfrac{1}{6}\|q\|_{L^6}^6-\tfrac1{4\pi}\|q\|_{L^4}^4\|q\|_{L^2}^2+\tfrac{1}{12\pi^2}\|q\|_{L^2}^6.
\end{align*}
\end{lemma}

\begin{proof}
To establish the expressions for $P(q)$ and $H(q)$ we first calculate
\begin{align*}
\mathcal{E}_1(q)&=\lan q,\Lc_qq\ran=\|q\|_{\dot{H}^{1/2}}^2\mp\lan |q|^2, C_+(|q|^2)\ran.
\end{align*}
Employing \eqref{Cid}, we may write
\begin{align*}
\lan |q|^2, C_+(|q|^2)\ran&=\lan |q|^2,|q|^2-C_-(|q|^2)+\tfrac1{2\pi}\mathcal{E}_0(q)\ran\\
&=\|q\|_{L^4}^4-\lan |q|^2,C_-(|q|^2)\ran+\tfrac1{2\pi}\mathcal{E}_0^2(q).
\end{align*}
As $\lan |q|^2,C_-(|q|^2)\ran=\overline{\lan |q|^2,C_+(|q|^2)\ran}=\lan |q|^2,C_+(|q|^2)\ran$, we find
\begin{equation}\label{eqn:norm_Cp_q^2}
\lan C_+ (|q|^2),C_+(|q|^2)\ran=\lan |q|^2,C_+(|q|^2)\ran=\tfrac12\|q\|_{L^4}^4+\tfrac1{4\pi}\|q\|_{L^2}^4.
\end{equation}
Thus,
\begin{align}\label{7:50}
\mathcal{E}_1(q)=\|q\|_{\dot{H}^{1/2}}^2\mp\tfrac12\|q\|_{L^4}^4\mp\tfrac1{4\pi}\|q\|_{L^2}^4
\end{align}
and the claimed expression for $P(q)$ follows.

As 
\begin{align}\label{5:08}
\mathcal{E}_2(q)&=\lan \Lc_qq,\Lc_qq\ran=\|q'\mp i qC_+(|q|^2)\|_{L^2}^2,
\end{align}
employing \eqref{7:50} we obtain
\begin{align*}
H(q)&=\tfrac12\mathcal{E}_2(q)\pm\tfrac{3}{4\pi}\mathcal{E}_0(q)\mathcal{E}_1(q)+\tfrac{1}{4\pi^2}\mathcal{E}_0(q)^3\\
&=\tfrac12\|q'\mp i qC_+(|q|^2)\|_{L^2}^2\pm \tfrac{3}{4\pi} \|q\|_{L^2}^2\|q\|_{\dot{H}^{1/2}}^2-\tfrac{3}{8\pi}\|q\|_{L^2}^2\|q\|_{L^4}^4+\tfrac{1}{16\pi^2}\|q\|_{L^2}^6.
\end{align*}

To obtain the detailed representation of $H(q)$ in the statement of this lemma, we have to further expand $\mathcal E_2(q)$.  We have
\begin{align}\label{eqn:E_2_expansion}
\mathcal{E}_2(q)&=\lan \Lc_qq,\Lc_qq\ran=\|q'\|_{L^2}^2\pm 2\Re\,\lan iq',q C_+(|q|^2)\ran+\|q C_+(|q|^2)\|_{L^2}^2.
\end{align}
First we study
\begin{align*}
2\Re\,\lan iq',q C_+(|q|^2)\ran&=2\Re\,\lan i\bar{q}q', C_+(|q|^2)\ran\\
&=2\Re\,\lan i\bigl[(|q|^2)'-q\overline{q'}\,\bigr], C_+(|q|^2)\ran.
\end{align*}
Averaging the two lines above, we obtain
\begin{align}\label{7:46}
2\Re\,\lan iq',q C_+(|q|^2)\ran&=\Re\,\lan i(|q|^2)', C_+(|q|^2)\ran+\Re\,\lan i(\bar{q}q'-q\overline{q'}), C_+(|q|^2)\ran.
\end{align}
The first term is
\begin{align}\label{7:47}
\Re\,\lan i(|q|^2)', C_+(|q|^2)\ran=-\| C_+(|q|^2)\|_{\dot{H}^{1/2}}^2.
\end{align}
For the second term, we use that $ C_+(|q|^2)=\overline{C_-(|q|^2)}$ to write
\begin{align*}
\Re\,\lan i(\bar{q}q'-q\overline{q'}), C_+(|q|^2)\ran
&=-2\Re\,\lan\Im(\bar{q}q'), C_+(|q|^2)\ran\\
%&=-2\lan\Im(\bar{q}q_x),\Re C_+(|q|^2)\ran\\
&=-2\Re\,\lan\Im(\bar{q}q'),C_-(|q|^2)\ran.
\end{align*}
Averaging again the last two lines and invoking \eqref{Cid} we obtain
\begin{align}\label{7:48}
\Re\, \lan i(\bar{q}q'-q\overline{q'}), C_+(|q|^2)\ran\notag
&=-\Re\,\lan \Im(\bar{q}q'),|q|^2+\tfrac1{2\pi}\mathcal{E}_0(q)\ran\notag\\
&=\Im\,\lan\bar{q}q',|q|^2\ran-\tfrac1{2\pi}\mathcal{E}_0(q)\Im\lan q,q'\ran\notag\\
&=\tfrac12 \Im\,\lan (q^2)',q^2\ran-\tfrac1{2\pi}\mathcal{E}_0(q)\|q\|_{\dot H^{1/2}}^2\notag\\
&=-\tfrac12\|q^2\|_{\dot{H}^{1/2}}^2-\tfrac1{2\pi}\|q\|_{L^2}^2\|q\|_{\dot{H}^{1/2}}^2.
\end{align}
Collecting \eqref{7:46}, \eqref{7:47}, and \eqref{7:48}, we obtain
\begin{align}\label{7:49}
2\Re\,\lan iq',q C_+(|q|^2)\ran=-\| C_+(|q|^2)\|_{\dot{H}^{1/2}}^2-\tfrac12\|q^2\|_{\dot{H}^{1/2}}^2-\tfrac1{2\pi}\|q\|_{L^2}^2\|q\|_{\dot{H}^{1/2}}^2.
\end{align}

Next, employing \eqref{Cid} we compute
\begin{align*}
\|q C_+(|q|^2)\|^2_{L^2}
&=\int_\T |q|^2 C_+(|q|^2)C_-(|q|^2)\,dx\\
&=\int_\T  \bigl[C_+(|q|^2)+C_-(|q|^2) -\tfrac1{2\pi} \mathcal{E}_0(q)\bigr]C_+(|q|^2)C_-(|q|^2)\,dx\\
&=\tfrac{1}{3}\int_\T\bigl[C_+(|q|^2)+C_-(|q|^2)\bigr]^3-\bigl[ C_+(|q|^2)\bigr]^3-\bigl[C_-(|q|^2)\bigr]^3\,dx\\
&\quad-\tfrac1{4\pi}\|q\|_{L^2}^2\|q\|_{L^4}^4-\tfrac1{8\pi^2}\|q\|_{L^2}^6,
\end{align*}
where we used \eqref{eqn:norm_Cp_q^2} in the final line.  Using again \eqref{Cid} we find
\begin{align*}
\int_\T\bigl[C_+(|q|^2)+C_-(|q|^2)\bigr]^3\,dx
&=\int_\T\bigl[|q|^2+\tfrac1{2\pi}\mathcal{E}_0(q)\bigr]^3\,dx\\
&=\|q\|_{L^6}^6+\tfrac3{2\pi}\|q\|_{L^4}^4\|q\|_{L^2}^2+\tfrac{1}{\pi^2}\|q\|_{L^2}^6.
\end{align*}
Using Fourier support considerations, we obtain
\begin{align*}
\int_\T\bigl[ C_+(|q|^2)\bigr]^3\,dx&=2\pi\bigl[\tfrac{1}{2\pi}\mathcal{E}_0(q)\bigr]^3=\tfrac{1}{4\pi^2} \|q\|_{L^2}^6
\end{align*}
and similarly,
\begin{align*}
\int_\T\bigl[ C_-(|q|^2)\bigr]^3\,dx&=2\pi\bigl[\tfrac{1}{2\pi}\mathcal{E}_0(q)\bigr]^3=\tfrac{1}{4\pi^2} \|q\|_{L^2}^6.
\end{align*}
We deduce that
\begin{align}
\|q C_+(|q|^2)\|^2_{L^2}
&=\tfrac{1}{3}\|q\|_{L^6}^6+\tfrac1{4\pi}\|q\|_{L^4}^4\|q\|_{L^2}^2+\tfrac{1}{24\pi^2}\|q\|_{L^2}^6.\label{eqn:L6_expression}
\end{align}

Combining \eqref{eqn:E_2_expansion} with \eqref{7:49} and \eqref{eqn:L6_expression}, we obtain 
\begin{align}
\mathcal{E}_2(q)&=\|q\|_{\dot{H}^{1}}^2\mp\bigl[\| C_+(|q|^2)\|_{\dot{H}^{1/2}}^2+\tfrac12\|q^2\|_{\dot{H}^{1/2}}^2+\tfrac1{2\pi}\|q\|_{L^2}^2\|q\|_{\dot{H}^{1/2}}^2\bigr]\nonumber\\
&\quad+\tfrac{1}{3}\|q\|_{L^6}^6+\tfrac1{4\pi}\|q\|_{L^4}^4\|q\|_{L^2}^2+\tfrac{1}{24\pi^2}\|q\|_{L^2}^6.\label{eqn:E2}
\end{align}

Inserting this into \eqref{5:08} we find
\begin{align*}
H(q)&= \tfrac12\|q\|_{\dot{H}^{1}}^2\mp\tfrac12\| C_+(|q|^2)\|_{\dot{H}^{1/2}}^2\mp\tfrac14\|q^2\|_{\dot{H}^{1/2}}^2\pm \tfrac1{2\pi}\|q\|_{L^2}^2\|q\|_{\dot{H}^{1/2}}^2\\
&\quad+\tfrac{1}{6}\|q\|_{L^6}^6-\tfrac1{4\pi}\|q\|_{L^4}^4\|q\|_{L^2}^2+\tfrac{1}{12\pi^2}\|q\|_{L^2}^6,
\end{align*}
as desired.
\end{proof}

\end{document}
